\section{When Can One Method Be Called Better?}
\label{sec:comparison}

\subsection{Primary estimand and paired comparisons}
The primary claim is \emph{higher verified completion probability under the
same execution policy}, not universal superiority. Let $T$ be the number of
independent eligible tasks and $R$ the balanced number of attempts per method
and task, both fixed before scoring. Define
\[
 \widehat p_m=\frac{1}{T}\sum_{t=1}^{T}\overline Y_{mt},\qquad
 \overline Y_{mt}=\frac{1}{R}\sum_{r=1}^{R}Y_{mtr},\qquad
 \widehat\Delta_{AB}=\frac{1}{T}\sum_{t=1}^{T}
 (\overline Y_{At}-\overline Y_{Bt}).
\]
Report $100\widehat\Delta_{AB}$ in percentage points. Each task has equal weight;
any deployment-weighted alternative must be fixed in advance. Reverse direction
is the same contrast: $\widehat\Delta_{BA}=-\widehat\Delta_{AB}$.
More attempts measure stochastic variation; they do not create more independent
tasks. Related tasks derived from one issue must remain in a shared analysis
cluster, or be removed under an eligibility rule fixed before scoring.

One frozen preparation artifact estimates performance conditional on that
artifact. Generalizing to the reliability of an authoring procedure requires
independent preparations and corresponding analysis of their variation.
Repeating task execution does not supply those observations.

A/B is the presently defined comparison. Once C is specified, a frozen three-arm
study has three unordered contrasts: A/B, A/C and B/C. The reverse directions
are not three additional tests. C cannot be designed using A/B's scored failures
and then compared on those same tasks as if held out. Use a new held-out cohort
for such a later comparison and rerun the reference workflows under its common
conditions.

\subsection{Prespecified analysis and error control}
Use a two-sided paired randomization test on the mean task-level difference,
keeping each task's entire vector of repeated outcomes together when swapping
method labels. Enumerate all swaps where feasible; otherwise use a frozen seed
and 100,000 random swaps, with the Monte Carlo correction
$(b+1)/(100{,}000+1)$, where $b$ counts absolute permuted statistics at least as
large as observed. Exact enumeration uses its exact tail fraction. Randomizing
method assignment to execution slots and preventing carryover are required
parts of the design. The test relies on exchangeability under the no-method-effect
null, not on treating every run as independent. With one binary observation per
task it reduces to the exact paired discordance test.

The adaptation of paired testing to this design follows the principles discussed
by \citet{dror2018}. Tasks are the resampling unit. If the eligible task structure
requires larger independent clusters, swap those clusters instead; freeze that
choice before scoring. If no credible exchangeable units exist, give descriptive
contrasts rather than an unsupported significance claim.

Set the primary family-wise error level to $0.05$. For $K$ planned pairwise
quality contrasts, apply Holm's step-down procedure \citep{holm1979}: order raw
$p$-values, compare the smallest with $0.05/K$, then successive values with
$0.05/(K-1)$ and so on, stopping at the first non-rejection. Report raw and
adjusted $p$-values. Here $K=1$ for A/B alone and $K=3$ for a prespecified
three-arm study. Additional confirmatory agent configurations enlarge the
registered family; do not choose the favorable model afterward.

For effect magnitude, report paired task-cluster bootstrap percentile intervals
at 95\%, resampling whole task vectors 100,000 times with a recorded seed.
These intervals are approximate, pointwise estimation intervals; they are not
simultaneous intervals and do not replace the multiplicity-adjusted test.
Small or nearly invariant samples can yield unreliable bootstrap intervals:
show the paired outcomes and flag that limitation rather than presenting false
precision. The power stage must validate this analysis for the intended sample.

\begin{table}[htbp]
\centering\small
\caption{Symmetric primary decision rule for any registered pair A/B.}
\label{tab:decisions}
\begin{tabularx}{\textwidth}{@{}Y Y@{}}
\toprule
Evidence & Permitted conclusion \\
\midrule
Positive completion difference and Holm-adjusted $p\leq0.05$ & A achieved higher verified completion than B under the tested policy. \\
\addlinespace
Negative completion difference and Holm-adjusted $p\leq0.05$ & B achieved higher verified completion than A under the tested policy. \\
\addlinespace
Neither condition holds, or data integrity is unresolved & Evidence is inconclusive about superiority; report the estimate and uncertainty. \\
\bottomrule
\end{tabularx}
\end{table}

Significance does not establish an adoption-worthy effect. A practically
meaningful improvement additionally needs a minimum worthwhile completion
difference justified before scoring from the intended workload and adoption
costs. Report whether the uncertainty excludes smaller effects. No company
requirement or practical-effect threshold has been supplied, so this document
does not invent one. Nonsignificance is not proof of equal performance
\citep{lakens2017}.

\subsection{Resource, efficiency and maintenance conclusions}
Report all-method resource totals, per-attempt distributions and paired changes,
including unsuccessful attempts. A supplementary operational measure is total
workflow resource per verified completion:
\[
 C_{m,j}=\frac{R_{m,j}(TR)}{\sum_{t,r}Y_{mtr}}.
\]
This allocates the declared preparation and upkeep horizon across all measured
attempts; it must not count only successful-run consumption. If no task succeeds,
the ratio has no finite useful estimate and must not be reported as zero.
Keep tokens, elapsed seconds and human minutes separate. Recompute ratios within
paired cluster resamples when intervals are estimable. Preparation uncertainty
requires independent preparation observations; resampling execution tasks alone
does not measure variation in context-authoring effort.

These are secondary estimates, not additional chances to declare the primary
winner. ``A used fewer tokens'' can be reported alongside lower or uncertain
correctness, provided both are explicit. ``A is more efficient at preserved
correctness'' requires a prespecified correctness non-inferiority margin,
adequate power, and a supported resource reduction. For a registered margin
$\delta_Q>0$, the lower confidence bound on $\Delta_{AB}$ must exceed
$-\delta_Q$. A resource reduction requires an upper bound below zero for
$R_{A,j}-R_{B,j}$, or below one for its ratio. These bounds must have error
control for the complete registered efficiency-claim family. Until a meaningful
margin and that family are frozen, efficiency claims remain descriptive.

For maintenance, compare observed checking/repair effort jointly with supported
claims, unresolved cases and coverage after shared changes. No-update behavior
is not a maintenance win when guidance becomes stale. A rigorous maintenance
superiority claim needs enough independent evolution units and prespecified
accuracy/coverage guardrails; otherwise report measured case-level trade-offs.

Do not force a total ranking. A can complete more tasks while B consumes fewer
resources. A primary correctness winner is not automatically the cheapest or
easiest to maintain. An overall organizational recommendation requires stated
priorities, cost accounting and acceptable trade-offs; no weighted score or
projected savings are assumed here.

\subsection{Sample planning and stopping}
Set the smallest completion difference worth detecting and the desired power
before final task selection. Use an independent pilot to estimate paired
outcome discordance, task heterogeneity and run-to-run variation. Simulate the
registered randomization test and Holm family over candidate task counts and
repetitions; select a design that reaches the declared power under the justified
planning assumptions. Preserve the simulation specification and sensitivity to
less favorable assumptions. Literature sample sizes are not our power analysis.

Independent tasks, rather than repeated attempts at one task, determine breadth
of evidence. If the feasible task pool cannot support the relevant effect, label
the study exploratory and report uncertainty. Freeze the task count, repetition
count, exclusions, analysis code and stopping rule before scoring. Do not keep
adding tasks until a significant result appears. A small pilot can establish
feasibility; it cannot guarantee a method-ranking claim.

\subsection{Research conclusion and remaining execution inputs}
The supported thesis direction is an incremental comparison against existing
project practice, with independently checked outcomes and full preparation and
upkeep accounting. The competing hypotheses remain live: preparation can remove
repeated discovery, or it can duplicate native guidance, consume attention and
increase upkeep. The study can distinguish these explanations in Memos; the
literature alone cannot select its winner.

The comparison definitions above are ready for protocol development. The
remaining inputs are concrete tasks and acceptance checks, B's preparation
procedure, an agent/model version, validated environment, execution limits,
independent-unit and sample/power choices, and maintenance changes. C needs its
own definition before inclusion. These are required experimental inputs, not
missing literature conclusions. There are no measured Memos outcomes yet.

A final result should name the tested workload, revision, agent and budget;
report the paired completion difference and uncertainty; apply the registered
pairwise decision; and state preparation, execution and maintenance trade-offs.
This supports ``better on verified completion in this setting'' when the data
justify it, and an explicit inconclusive result when they do not.
